\documentclass[10pt,twoside]{article}
\newcommand{\BedrockArchitectureName}{Bedrock}
\newcommand{\BedrockDocumentTitle}{Bedrock Target Intrinsics}
\newcommand{\BedrockRunningTitle}{TARGET INTRINSICS}
\newcommand{\BedrockFooterTitle}{BEDROCK ARCHITECTURE}
\usepackage{bedrock-reference}
\begin{document}
\BedrockMakeTitlePage{Bedrock Target Intrinsics}{C Compiler Builtins and Header Interfaces}{Draft}
\BedrockDocumentHasFiguresfalse
\BedrockFrontMatter

\clearpage
\section{Scope and Status}\label{page:scope-and-status}
This document defines the compiler-owned Bedrock C target interface: builtin
names, implementation-reserved header wrappers, source signatures, lowering
requirements, availability and privilege constraints, and compiler-visible
effects. It is the contract shared by a Bedrock-aware C frontend, optimizer,
backend, and the target headers.

Architectural instruction behavior remains defined by the Bedrock ISA
Specification. The baseline calling convention is defined by the Bedrock C
ABI. This document specifies how source programs reach target facilities; it
does not duplicate their architectural or language semantics.

\BedrockField{Interface name:}{\texttt{bedrock-target-intrinsics}}
\BedrockField{Source language:}{ISO C with implementation-reserved target interfaces}

\section{Compiler Interface Model}\label{page:compiler-interface-model}

\subsection{Names and Ownership}
The compiler builtin namespace begins with
\texttt{\_\_builtin\_bedrock\_}. Each \textbf{Name} entry in this document
denotes that prefix followed by the displayed name. For example,
\texttt{cpuid} denotes
\texttt{\_\_builtin\_bedrock\_cpuid}. The ordinary header wrapper uses the
\texttt{\_\_bedrock\_} prefix with the same name unless its entry specifies a
macro spelling.

All of these identifiers are implementation-reserved. A builtin is a compiler
primitive, not an external function symbol. The target headers own wrapper
declarations and umbrella include topology; this document owns their
ABI-visible meaning.

\subsection{Header Topology and Exposure}
\texttt{<bedrockintrin.h>} is the public umbrella. It includes the core, memory,
integer, floating-point, and vector families. \texttt{<bedrocksystemintrin.h>} is the system
umbrella for system-register, cache, and MMU facilities. Any
family header may also be included directly.

(:intrinsic-groups:cache,core,fpu,integer,memory,mmu,sysreg,vector:)

\subsection{Common Compilation Rules}
Family wrappers are static inline functions unless an encoded immediate or
type operand requires a macro. Every builtin listed by an enabled profile must
be recognized by the compiler. If a required optional feature is disabled,
compilation fails rather than emitting an unresolved call.

An encoded immediate must be an integer constant expression in the stated
range. Header inclusion neither suppresses nor emulates architectural
privilege checks. An operation is volatile only when its entry states an
externally observable architectural effect. Compiler memory effects and
hardware ordering are exactly those stated by the entry; one does not imply
the other.

The \textbf{C interface} column uses
\texttt{\detokenize{result(arguments)}} shorthand. Exact declarations are in the
named header; \texttt{u8}, \texttt{u16}, \texttt{u32}, and \texttt{u64}
abbreviate the corresponding \texttt{uintN\_t} types, and \texttt{result}
denotes \texttt{\_\_bedrock\_query\_result\_t}. A clobber means a compiler
memory clobber, while a named fence also has the hardware behavior defined by
the ISA.

\section{Public Target Interface}\label{page:public-target-interface}

\subsection{Core Builtins}
\BedrockField{Header:}{\texttt{<bedrockcoreintrin.h>}}
(:intrinsic-group:core:)

\subsection{Memory Builtins}
\BedrockField{Header:}{\texttt{<bedrockmemoryintrin.h>}}
(:intrinsic-group:memory:)

\subsection{Integer Builtins}
\BedrockField{Header:}{\texttt{<bedrockintegerintrin.h>}}
(:intrinsic-group:integer:)

\subsection{Floating-Point Builtins}
\BedrockField{Header:}{\texttt{<bedrockfpuintrin.h>}}
(:intrinsic-group:fpu:)

\subsection{Vector Builtins}
\BedrockField{Header:}{\texttt{<bedrockvectorintrin.h>}}
(:intrinsic-group:vector:)

\subsection{Approximate Transcendental Lowering Policy}
A strict language-library operation shall not be lowered to an FPTRANSA
instruction. A compiler may select FPTRANSA only when the source uses a
separately defined approximate interface or an active compiler mode permits
approximate mathematics. The selected instruction remains subject to its ISA
accuracy, special-value, and floating-point exception-condition contract.

\section{System Target Interface}\label{page:system-target-interface}
The compiler emits the requested architectural operation; runtime privilege and policy checks
remain active.

\subsection{System-Register Builtins}
\BedrockField{Header:}{\texttt{<bedrocksysregintrin.h>}}
(:intrinsic-group:sysreg:)

The system-register header defines \texttt{\_\_bedrock\_control\_register\_t} constants for every published
control-register selector. The constants are names for the architectural selector values; they do not weaken the
privilege, reserved-bit, or per-register validation performed by RDCR and WRCR.

It also defines the pure value macros
\texttt{\_\allowbreak{}\_\allowbreak{}BEDROCK\_\allowbreak{}SEGMENT\_\allowbreak{}IMAGE},
\texttt{\_\allowbreak{}\_\allowbreak{}BEDROCK\_\allowbreak{}SEGMENT\_\allowbreak{}IMAGE\_\allowbreak{}FOR\_\allowbreak{}BASE}, and
\texttt{\_\allowbreak{}\_\allowbreak{}BEDROCK\_\allowbreak{}SEGMENT\_\allowbreak{}DISABLED}. They construct 64-bit segment
images without reading or writing architectural state.

\texttt{\_\_BEDROCK\_SEGMENT\_IMAGE(base\_page,exponent,mantissa,
bounds\_only)} places the low 52, 5, and 6 bits of the first three operands in
their architectural fields and normalizes \texttt{bounds\_only} to zero or
one. Each operand is evaluated once. A valid enabled image uses an exponent
from 0 through 31 and a mantissa from 1 through 63. The macro only packs
fields; it does not validate the resulting base, span, limit, or address
canonicality.

\texttt{\_\_BEDROCK\_SEGMENT\_IMAGE\_FOR\_BASE(base,...)} encodes
$\lfloor\texttt{(uint64\_t)base}/4096\rfloor$ in the base-page field. It
discards the low 12 bits, so a misaligned operand names the containing page and
does not preserve the original byte address. The macro imposes no alignment
precondition and does not adjust a later offset to compensate for the discarded
bits. \texttt{\_\_BEDROCK\_SEGMENT\_DISABLED} is the canonical all-zero
disabled image.

\subsection{Cache-Management Builtins}
\BedrockField{Header:}{\texttt{<bedrockcacheintrin.h>}}
All cache range operations evaluate the address once and obtain the maintenance granule from CPUID
\texttt{CACHE\_TOPOLOGY}. A zero length emits no cache-maintenance instruction. A nonzero, nonwrapping byte range
selects every maintenance block that intersects the range and lowers to one block instruction per selected block.
Each instruction translates its own block address. If a later instruction faults, blocks completed by earlier
instructions remain completed.

For a nonzero length, the mathematical half-open range
$[\mathrm{address},\mathrm{address}+\mathrm{length})$ shall lie within the
64-bit address space. Otherwise the behavior of the call is undefined and no
architectural cache-maintenance sequence is required. An implementation may
diagnose a provably wrapping constant range. A valid range end is calculated
mathematically, not modulo $2^{64}$.

(:intrinsic-group:cache:)

\subsection{MMU Builtins}
\BedrockField{Header:}{\texttt{<bedrockmmuintrin.h>}}
(:intrinsic-group:mmu:)

\section{Shared Header Types}\label{page:shared-header-types}
(:intrinsic-types:base.types.core.performance_counter,base.types.mmu.query_result,base.types.sysreg.control_register,base.types.sysreg.segment_register:)

\end{document}
